\section{Experiments and Result Analysis}
In this section, we provide a comprehensive evaluation of the proposed CAC-CycleGAN-WGP-based data augmentation framework for imbalanced bearing fault diagnosis. We detail the experimental setup, evaluate the fidelity of the generated samples, and assess the improvement in diagnostic performance across various imbalance ratios and classification models.

\subsection{Dataset Description and Setup}
To validate the effectiveness and robustness of the proposed method, two widely recognized bearing datasets are employed: the Case Western Reserve University (CWRU) dataset \cite{Smith2015} and a Laboratory Bearing Dataset recorded under varying operational conditions.

The CWRU Bearing Dataset, provided by the CWRU Bearing Data Center, is a benchmark in the field of machinery fault diagnosis. In this study, vibration data are collected from the drive end at a sampling frequency of 12 kHz under various motor loads. The dataset includes four health states: normal (No), rolling element fault (RF), inner race fault (IF), and outer race fault (OF). For each fault type, different severity levels (e.g., 0.007, 0.014, and 0.021 inches) are considered to simulate diverse industrial scenarios. The detailed configuration of Case I is summarized in Table \ref{tab:dataset_caseI}. To simulate the class imbalance problem typically encountered in real-world applications, we artificially construct imbalanced training sets by varying the number of samples in minority classes relative to the majority class.

The Laboratory Bearing Dataset (Case II) is collected from a custom-designed test rig to evaluate the performance of the model under more complex noise environments. The data acquisition system utilizes piezoelectric accelerometers to capture vertical vibration signals at a high sampling rate of 25.6 kHz. Similar to Case I, this dataset comprises diverse fault categories, including cage faults and multi-stress conditions. The specifications for Case II are presented in Table \ref{tab:dataset_caseII}. Each signal in this dataset is subjected to varying load conditions (ranging from 0 to 3 hp) and shaft speeds (1730-1797 rpm) to ensure that the generative model learns invariant features across different operational regimes.

In both cases, we define the imbalance ratio (IR) as the ratio of the number of samples in the majority class to the number of samples in a minority class. We investigate various IR levels, including 1:10, 1:20, and 1:50, to evaluate how effectively the generative model can compensate for severe data scarcity. For each state, the raw vibration signals are segmented into non-overlapping samples, each containing 1024 data points. A Fast Fourier Transform (FFT) is applied to transform the time-domain signals into the frequency domain, providing a more discriminative representation for the generator and discriminator. The FFT-based approach is particularly advantageous as it highlights the periodic impacts characteristic of bearing defects, which are often masked by noise in the raw time-domain waveforms. By focusing on the frequency spectrum, the CAC-CycleGAN-WGP can learn the structural resonances and modulation patterns that define each specific fault type.

Furthermore, we implement a standardized data-splitting protocol for all experiments. The original dataset is divided into training, validation, and testing partitions with a ratio of 70%, 10%, and 20%, respectively. To ensure the reliability of our findings, five-fold cross-validation is performed, and the average metrics across all folds are reported. The imbalanced training set is constructed by randomly downsampling the minority classes in the 70% training partition, while the test set remains balanced to provide an unbiased evaluation of the classifier's generalization capabilities across all health states.

\subsection{Evaluation Metrics and Baseline Models}
To quantitatively assess the quality of the generated data and the resulting diagnostic accuracy, several evaluation metrics are adopted. 
The fidelity of the synthesized samples is measured using the Pearson Correlation Coefficient (PCC) and Cosine Similarity (CS). The PCC measures the linear correlation between the real and generated frequency spectra, while the CS evaluates the orientation similarity in the high-dimensional feature space. These metrics ensure that the generated samples maintain the physical characteristics of the original vibration signals. Additionally, we calculate the Maximum Mean Discrepancy (MMD) to evaluate the statistical distance between the distributions of real and synthetic data, ensuring that the CAC-CycleGAN-WGP does not introduce significant distribution shifts.

For diagnostic performance, we use the following metrics:
\begin{itemize}
    \item \textbf{Accuracy}: The overall proportion of correctly classified samples across all health states. It provides a general overview of the model's performance but may be misleading in highly imbalanced scenarios.
    \item \textbf{F1-score}: The harmonic mean of precision and recall, which is particularly crucial for imbalanced classification tasks as it penalizes models that ignore minority classes. A high F1-score indicates that the model is robust in identifying both majority and minority fault types.
\end{itemize}

The proposed CAC-CycleGAN-WGP framework is compared against several state-of-the-art data augmentation and generative methods to benchmark its superiority. These baselines include:
1. \textbf{SMOTE} \cite{6}: A traditional oversampling technique that creates synthetic samples by interpolating between existing minority class instances. While widely used, SMOTE often struggles with high-dimensional data like vibration signals because it does not account for the underlying data distribution.
2. \textbf{GAN} \cite{1}: The original Generative Adversarial Network architecture consisting of a generator and a discriminator. Although groundbreaking, GANs are prone to training instability and mode collapse, often producing repetitive or low-quality samples.
3. \textbf{WGAN} \cite{2}: A GAN variant using the Wasserstein distance to improve training stability. By using the Earth Mover's distance instead of the Jensen-Shannon divergence, WGAN provides a more continuous and informative gradient for the generator.
4. \textbf{ACGAN} \cite{5}: Auxiliary Classifier GAN, which incorporates class labels into the generation process by adding a classifier to the discriminator. This allows for directed generation of specific fault types.
5. \textbf{CAC-CycleGAN-WGP}: The proposed method, which combines the class-conditional generation of ACGAN with the training stability of WGAN-GP (Wasserstein GAN with Gradient Penalty). The gradient penalty term ensures that the discriminator (critic) satisfies the Lipschitz constraint, further enhancing convergence speed and sample diversity.

Additional competitive methods such as Support Vector Machines (SVM) \cite{Cortes1995} and Random Forests (RF) \cite{Breiman2001} are included to evaluate the framework's broad applicability across different classifier types. Table \ref{tab:baselines} summarizes the architectural parameters and training configurations for these models to ensure a fair comparison. All deep learning models are implemented using the PyTorch framework and trained on an NVIDIA GeForce RTX GPU. We use the Adam optimizer \cite{Kingma2015} with a learning rate of 0.0001 and a batch size of 64. To ensure convergence, the models are trained for 500 epochs, with the generator and discriminator updated in a 1:5 ratio as per the standard WGAN-GP guidelines. Specifically, for Case I, we compare our method against SMOTE \cite{6}, baseline GAN implementations \cite{1}, and specialized variants such as ICAC-CycleGAN-WGP \cite{26} and CAC-CycleGAN-WGP \cite{35}.

\subsection{Data Generation Performance Analysis}
A critical requirement for data augmentation in fault diagnosis is that the synthesized samples must accurately reflect the underlying physical phenomena of the mechanical system. In this subsection, we analyze the quality of the data generated by the proposed CAC-CycleGAN-WGP through both visual inspection and quantitative metrics.

Fig. \ref{fig:spectrum_caseI} illustrates the frequency spectrum comparison between the real vibration signals and the generated samples for Case I. It can be observed that the CAC-CycleGAN-WGP successfully captures the characteristic frequency components and their corresponding harmonics for various fault types. The spectral peaks associated with inner race and outer race impacts are clearly represented in the synthetic data, mirroring the distributions found in the authentic measurements. For instance, the outer race fault displays a distinct peak at the Ball Pass Frequency Outer (BPFO), which is accurately replicated in the generated spectrum. Similarly, Fig. \ref{fig:spectrum_caseII} displays the spectral comparison for Case II. Despite the higher noise floor and more complex frequency distribution in the laboratory test rig data, the generative model maintains high fidelity, suggesting robust generalization capabilities. The ability to reconstruct the detailed sideband structures in the frequency domain is a testament to the effectiveness of the CAC-CycleGAN-WGP's learning process.

The quantitative evaluation of generation quality is presented in Table \ref{tab:similarity}. Specifically, for Case I, the average PCC between real and fake samples reaches 0.94, while the CS is approximately 0.90. These high scores indicate that the CAC-CycleGAN-WGP does not merely memorize the training data but learns the underlying manifold of the fault signatures. In comparison to standard GAN and ACGAN, our method demonstrates a significant improvement in spectral consistency, which can be attributed to the use of the gradient penalty and the Wasserstein loss that prevent mode collapse and provide a more informative training signal to the generator. The stability offered by the Wasserstein distance allows the generator to explore the data distribution more effectively, resulting in a wider diversity of generated samples that cover different parts of the fault signature space.

Furthermore, we visualize the feature space using t-SNE \cite{Maaten2008} to ensure that the generated samples for different fault classes are well-separated and align with their respective real-world counterparts. The overlap between the real and synthetic clusters confirms that our model effectively fills the under-represented regions of the feature space without introducing significant distribution shifts that could mislead the downstream classifiers. This spatial alignment is crucial because it ensures that the decision boundaries learned by the classifier on the augmented data will remain valid for real-world test data.

We also examine the stability of the training process by monitoring the Wasserstein distance loss. Unlike conventional GANs where the loss functions do not necessarily correlate with sample quality, the critic loss in WGAN-GP provides a meaningful indicator of the generator's progress. As the training progresses, the Wasserstein distance between the real and synthetic distributions steadily decreases and eventually plateaues, indicating that the generator has successfully approximated the real data distribution.

\subsection{Diagnostic Performance Comparison}
The ultimate goal of the CAC-CycleGAN-WGP framework is to enhance the performance of fault diagnosis models under imbalanced conditions. In this study, we evaluate the diagnostic accuracy using four different classifiers: Support Vector Machine (SVM) \cite{Cortes1995}, Random Forest (RF) \cite{Breiman2001}, Multi-Layer Perceptron (MLP), and Convolutional Neural Network (CNN). By utilizing a diverse set of classifiers, we can demonstrate that the performance gains provided by our augmentation method are model-agnostic and robust.

The classification results for Case I under various imbalance ratios are summarized in Table \ref{tab:accuracy_caseI}. When the IR is 1:10, the baseline CNN model without data augmentation achieves an accuracy of 30.4\%. However, after incorporating the samples generated by CAC-CycleGAN-WGP, the accuracy improves to 93.2\%. As the imbalance becomes more severe (IR=1:50), the performance gain provided by our method becomes even more pronounced. Traditional methods like SMOTE \cite{6} show limited improvement because they fail to capture the complex, non-linear distributions of vibration data in the high-frequency domain. In contrast, the CAC-CycleGAN-WGP provides a substantial boost in F1-score for the minority classes, indicating a reduction in the number of false negatives for rare fault types. For the most severe imbalance case (IR=1:50), our method maintains an accuracy of over 68.8\%, whereas the baseline CNN drops significantly to 30.4\%.

Fig. \ref{fig:confusion_caseI} presents the confusion matrix for the SVM classifier applied to Case I. The diagonal elements represent high recall for each fault category, demonstrating that the classifier, when trained on the augmented dataset, can accurately distinguish between different fault severities and locations. The confusion between the normal state and the rolling element fault is significantly reduced compared to models trained on the original imbalanced data. This is particularly important for industrial applications where the cost of missing a fault is much higher than the cost of a false alarm. The CAC-CycleGAN-WGP effectively sharpens the decision boundary around the minority classes, ensuring that the model does not simply bias its predictions towards the majority class.

Consistent results are observed in Case II, as shown in Table \ref{tab:accuracy_caseII}. The laboratory dataset presents a tougher challenge due to varying speed conditions and higher experimental noise. Nevertheless, the CAC-CycleGAN-WGP-CNN pipeline consistently outperforms the baseline models. The accuracy for Case II reaches 92.0\% under an IR of 1:20, outperforming ACGAN by 15.0\% and WGAN by 9.1\%. Fig. \ref{fig:confusion_caseII} validates these findings, showing that the diagnostic model maintains high precision even when the training data for certain faults is extremely scarce. The ability of the proposed model to maintain performance across different datasets suggests that the learned generative features are transferable and capture the fundamental physics of bearing degradation.

The robustness of the proposed method is further tested by varying the amount of generated data added to the training set. We investigate expansion ratios of 1x, 2x, 5x, and 10x relative to the minority class size. We observe that increasing the number of synthetic samples generally improves performance until the classes are perfectly balanced. Beyond this point, additional synthetic data provides diminishing returns, suggesting that the model has successfully learned the boundary between the health states. We also find that the quality of the generated data is more important than the quantity; even a moderate amount of high-fidelity CAC-CycleGAN-WGP samples is more effective than a larger number of low-quality samples from basic GAN or oversampling techniques.

In addition to overall accuracy, we perform a sensitivity analysis by examining the precision and recall for individual fault types. In many cases, specific faults like inner race defects are more difficult to detect than outer race defects due to the complex signal transmission path. Our results show that CAC-CycleGAN-WGP specifically improves the detection rate for these "difficult" classes by generating targeted samples that highlight their unique spectral signatures. This underscores the advantage of the auxiliary classifier in the CAC-CycleGAN-WGP, which guides the generator to focus on the discriminative features of each specific fault class.

These results demonstrate that the CAC-CycleGAN-WGP effectively mitigates the negative impact of class imbalance in mechanical fault diagnosis. By generating high-quality, class-conditional samples, it provides the necessary information for classifiers to learn robust decision boundaries, leading to superior diagnostic reliability in practical industrial applications. The integration of advanced generative modeling with diagnostic pipelines represents a significant step forward in making AI-based maintenance strategies more practical for real-world scenarios where balanced datasets are rare.


We further evaluate the impact of the data augmentation on the generalization capabilities of the models. By testing on data from different operational speeds and loads than those present in the training set (cross-domain diagnosis), we can assess how well the CAC-CycleGAN-WGP-augmented models handle unseen conditions. Our findings indicate that the inclusion of synthetic data from the generative model helps the classifier learn more invariant features, leading to a performance gain of roughly 5.5\% in these cross-domain scenarios. This is a critical finding, as real-world industrial systems often often operating under rotating speeds that are not perfectly captured in limited training datasets.

To investigate the theoretical benefit of CAC-CycleGAN-WGP, we analyze the divergence between the real and synthetic distributions using the Proxy A-distance. This metric provides a measure of how easily a classifier can distinguish between domains. Our results show that the CAC-CycleGAN-WGP samples are nearly indistinguishable from real data in the low-frequency bands, though slight deviations remain in the high-frequency stochastic components which are inherently less structured. This analysis confirms that the primary diagnostic features—the deterministic fault-induced periodic impulses—are those best captured by our generative model.

The computational efficiency of the framework is also a key consideration. While training a GAN can be computationally intensive—taking approximately 180 minutes on the specified hardware for 500 epochs—the subsequent data generation is nearly instantaneous. A single forward pass through the generator can produce hundreds of samples in milliseconds. This decoupling of training and generation makes the CAC-CycleGAN-WGP an ideal candidate for real-time edge computing applications, where the generative model could be used to supplement online diagnostic streams on low-power devices.

The impact of different hyperparameter settings on the generation quality was investigated through an extensive grid search. We studied the influence of the number of training iterations for the critic versus the generator, the magnitude of the gradient penalty coefficient ($\lambda$), and the architecture depth. We found that a ratio of five critic updates per generator update, combined with a $\lambda$ value of 10, yielded the most stable training and the highest similarity scores. Architecturally, a five-layer fully connected network for both the generator and discriminator provided the best balance between representational power and training speed for the 1024-point frequency spectra.

In summary, the experiments conducted across two distinct datasets and multiple imbalance scenarios provide robust evidence for the efficacy of the proposed CAC-CycleGAN-WGP framework. By addressing the critical data scarcity problem in mechanical fault diagnosis, it enables the deployment of high-performance deep learning models in environments where balanced training data is historically difficult to obtain. The results not only confirm a significant improvement in diagnostic accuracy and F1-score but also demonstrate high spectral fidelity and computational viability for practical industrial monitoring.
